\section{Related Work}

\subsection{Low-Rank Attention}

\textbf{SOA Fused Attention} and \textbf{Adaptive SVD Attention} factorize Q/K/V 
(and possibly O) via truncated SVD. \textbf{A3} and \textbf{AASVD} use 
activation-weighted SVD requiring calibration. VSpec++ differs by 
combining \textbf{spectral modulation} with \textbf{token routing}.

\subsection{Token Sparse Attention}

\textbf{H2O} (Zhang et al., 2023) evicts KV tokens by accumulated attention. 
\textbf{StreamingLLM} (Xiao et al., 2024) preserves attention sinks. \textbf{TCA} 
(companion work) uses a learnable saliency net with \texttt{[CLS]} preservation. 
VSpec++ uses a similar router, combined with the other gates.

\subsection{Head Sparsity}

\textbf{Squeezing Heads} (companion) distills heads. \textbf{Selective Attention} 
(Lee et al., 2019) prunes heads. \textbf{HeadGate} in VSpec++ provides per-example 
soft weighting rather than hard pruning.

\subsection{Variational Methods in Attention}

\textbf{Variational Dropout} (Kingma et al., 2015; Molchanov et al., 2017) 
learns per-weight dropout rates via reparameterization. \textbf{VIB} (Alemi 
et al., 2017) uses variational bounds for representation learning. 
\textbf{Variational Attention} (Bahuleyan et al., 2018) uses a latent variable 
for attention. Our variational spectral gate is closest to Variational 
Dropout but applied to singular value importances.

\subsection{Reparameterization Trick}

\textbf{Kingma \& Welling (2013)} introduced the reparameterization trick for 
training VAEs. \textbf{Rezende et al. (2014)} developed similar ideas 
independently. We use the standard Gaussian reparameterization:

\[
z = \mu + \sigma \odot \epsilon, \quad \epsilon \sim \mathcal{N}(0, I)
\tag{2.1}
\]

\subsection{Position of Our Work}

\begin{table}[H]
\centering
\small
\begin{tabular}{llllll}
\toprule
\textbf{Method} & \textbf{SVD} & \textbf{Token route} & \textbf{Spectral gate} & \textbf{Head gate} & \textbf{Variational} \\
\midrule
SOA Fused & Yes & No & No & No & No \\
TCA & No & Yes & No & No & No \\
Squeezing Heads & No & No & No & Yes & No \\
Variational Dropout & No & No & No & No & Yes \\
\textbf{VSpec++} & \textbf{Yes} & \textbf{Yes} & \textbf{Yes} & \textbf{Yes} & \textbf{Yes} \\
\bottomrule
\end{tabular}
\end{table}

\textbf{Key distinction.} VSpec++ is the first to combine all four techniques 
in a single attention block.



\subsection{Comparison Table}

\begin{table}[H]
\centering
\small
\begin{tabular}{lllll}
\toprule
\textbf{Method} & \textbf{Params} & \textbf{Token} & \textbf{Spectral} & \textbf{Variational} \\
\midrule
MHSA & $4d^2$ & Full & -- & -- \\
SOA Fused & $2d^2 + 2rd$ & Full & -- & -- \\
TCA & $4d^2$ & Top-$k$ & -- & -- \\
\textbf{VSpec++} & $2rd + \text{gates}$ & \textbf{Top-$k$} & \textbf{Yes} & \textbf{Yes} \\
\bottomrule
\end{tabular}
\end{table}

